% Options for packages loaded elsewhere
\PassOptionsToPackage{unicode}{hyperref}
\PassOptionsToPackage{hyphens}{url}
\documentclass[
]{article}
\usepackage{xcolor}
\usepackage{amsmath,amssymb}
\setcounter{secnumdepth}{-\maxdimen} % remove section numbering
\usepackage{iftex}
\ifPDFTeX
  \usepackage[T1]{fontenc}
  \usepackage[utf8]{inputenc}
  \usepackage{textcomp} % provide euro and other symbols
\else % if luatex or xetex
  \usepackage{unicode-math} % this also loads fontspec
  \defaultfontfeatures{Scale=MatchLowercase}
  \defaultfontfeatures[\rmfamily]{Ligatures=TeX,Scale=1}
\fi
\usepackage{lmodern}
\ifPDFTeX\else
  % xetex/luatex font selection
\fi
% Use upquote if available, for straight quotes in verbatim environments
\IfFileExists{upquote.sty}{\usepackage{upquote}}{}
\IfFileExists{microtype.sty}{% use microtype if available
  \usepackage[]{microtype}
  \UseMicrotypeSet[protrusion]{basicmath} % disable protrusion for tt fonts
}{}
\makeatletter
\@ifundefined{KOMAClassName}{% if non-KOMA class
  \IfFileExists{parskip.sty}{%
    \usepackage{parskip}
  }{% else
    \setlength{\parindent}{0pt}
    \setlength{\parskip}{6pt plus 2pt minus 1pt}}
}{% if KOMA class
  \KOMAoptions{parskip=half}}
\makeatother
\setlength{\emergencystretch}{3em} % prevent overfull lines
\providecommand{\tightlist}{%
  \setlength{\itemsep}{0pt}\setlength{\parskip}{0pt}}
\usepackage{bookmark}
\IfFileExists{xurl.sty}{\usepackage{xurl}}{} % add URL line breaks if available
\urlstyle{same}
\hypersetup{
  hidelinks,
  pdfcreator={LaTeX via pandoc}}

\author{}
\date{}

\begin{document}

\section{WILEY}\label{wiley}

\textit{[Image omitted: images/7889807a529ef4b9ac28cb518161ebeca896c50001ef02a52a8defdb9e5c5333.jpg]}

text\_image

ROYAL STATISTICAL SOCIETY FOUNDED 1834

Maximum Likelihood Estimation of Observer Error-Rates Using the EM
Algorithm

Author(s): A. P. Dawid and A. M. Skene

Reviewed work(s):

Source: Journal of the Royal Statistical Society. Series C (Applied
Statistics), Vol. 28, No.~1 (1979), pp.~20-28

Published by: Wiley for the Royal Statistical Society

Stable URL: http://www.jstor.org/stable/2346806

Accessed: 17/01/2013 01:52

Your use of the JSTOR archive indicates your acceptance of the Terms \&
Conditions of Use, available at
http://www.jstor.org/page/info/about/policies/terms.jsp

JSTOR is a not-for-profit service that helps scholars, researchers, and
students discover, use, and build upon a wide range of content in a
trusted digital archive. We use information technology and tools to
increase productivity and facilitate new forms of scholarship. For more
information about JSTOR, please contact support@jstor.org.

\textit{[Image omitted: images/4cc6a4132299aeebd61abf270aebf37766c31f178e15d02b1626758b6b542188.jpg]}

Wiley and Royal Statistical Society are collaborating with JSTOR to
digitize, preserve and extend access to Journal of the Royal Statistical
Society. Series C (Applied Statistics).

http://www.jstor.org

\section{Maximum Likelihood Estimation of Observer Error-rates using the
EM
Algorithm}\label{maximum-likelihood-estimation-of-observer-error-rates-using-the-em-algorithm}

By A. P. DAWID† and A. M. SKENE‡

University College London

{[}Received October 1977. Revised May 1978{]}

\section{SUMMARY}\label{summary}

In compiling a patient record many facets are subject to errors of
measurement. A model is presented which allows individual error-rates to
be estimated for polytomous facets even when the patient's ``true''
response is not available. The EM algorithm is shown to provide a slow
but sure way of obtaining maximum likelihood estimates of the parameters
of interest. Some preliminary experience is reported and the limitations
of the method are described.

Keywords: EM ALGORITHM; OBSERVER VARIATION; LATENT CLASS MODEL; MEDICAL
EXAMPLE

\section{1. INTRODUCTION}\label{introduction}

WHEN a patient's history is taken by different clinicians, different
replies may be obtained to the same question. This may occur for a
number of reasons; perhaps slightly different wording is used in each
case, or perhaps the question is one the patient finds difficult to
answer satisfactorily and so changes his reply from time to time.
Similarly, in classifying a facet (sign or symptom) for type, severity,
extent or duration, the patient and the clinicians may have different
interpretations of the underlying scale of measurement. Such facets are
said to be subject to observer error in that the response recorded may
not be the ``true'' response as defined by some standard description of
the facet or as implied by a consensus of medical opinion. The
consequences of observer error are investigated by Good and Card (1971),
where it is shown that a fairly low rate of error can lead to a
considerable loss of diagnostic information.

In this paper we consider the problem of measuring observer error when
the facet being recorded can take one of a finite number of values coded
1,2,\ldots,J. Two objectives can be distinguished. In some situations it
is of interest to measure the degree of observer agreement. Landis and
Koch (1977) give a general statistical methodology for the analysis of
multivariate categorical data involving agreement among more than two
observers. In fact the measurement of inter- and intra-observer
agreement has received wide attention. Landis and Koch (1975a, b) give a
review.

There are situations, however, where the performance of individual
observers is more relevant. Let \(\pi_{jl}^{(k)}\) be the probability
that an observer, \(k\) , will record value \(l\) given \(j\) is the
true response. The probabilities
\(\pi_{jl}^{(k)}, j = 1, \dots, J, l = 1, \dots, J\) are called the
individual error-rates for the \(k\) th observer, although this set
includes \(\pi_{jj}^{(k)}, j = 1, \dots, J\) , which are the
probabilities that the observer records the true response in each of the
\(J\) possible cases. Note that the error-rates are conditional
probabilities where

\[
\sum_ {l = 1} ^ {J} \pi_ {j l} ^ {(k)} = 1 \quad \text {for each} j \text {and} k,
\]

† Now at Maths Dept, The City University, Northampton Square, London
EC1V 0HB.

‡ Now at Maths Dept, University Park, Nottingham NG7 2RD.

and these quantities are mathematically independent of the marginal
distribution of the true response.

Individual error-rates can be used to advantage in several situations:

\begin{enumerate}
\def\labelenumi{(\roman{enumi})}
\tightlist
\item
  Through a knowledge of his own performance, an observer is able to
  recognize those facets where he makes the most frequent
  misclassifications and so reduce the number of errors in the patient's
  record. For example, he may reword a question or devote more time to
  the measurement of a particular sign. In those circumstances where a
  facet is recorded on an ordinal scale, error-rates serve to indicate
  how one observer's interpretation of the scale may differ from the
  general consensus, and thus show the way the observer should modify
  future assessments.\\
\item
  In the development of a data-base for purposes of diagnosis it is
  desirable to minimize errors of measurement. Knowledge of individual
  error-rates allows a contributor to that database to be monitored.\\
\item
  In some situations it might be recognized that a particular facet is
  subject to considerable observer error even when this is elicited by
  the most experienced clinicians. Nevertheless, it is highly desirable
  that this facet be accurately recorded and so, if several observers
  participate, a consensus judgement can be obtained. This judgement may
  be a simple majority opinion or a weighted consensus where the weights
  are functions of the individual error rates. In the latter case, each
  observer's contribution to the consensus is determined by his previous
  performance in eliciting that facet.\\
\item
  There has been recent interest in delegating certain tasks such as
  history-taking, which have traditionally been carried out by doctors,
  to ancillary personnel, or even to a computer terminal. It is
  important to know whether any loss (or gain) of accuracy is likely to
  ensue.
\end{enumerate}

Henceforth all discussion is with reference to a single facet. When the
true response can be obtained by some independent means, e.g.~passage of
time, X-ray or eventual operation, then sensible estimates of individual
error-rates can be calculated very easily. For example, one possible
estimator is

\[
\hat {\pi} _ {j l} ^ {(k)} = \frac {\text { number   of   times   observer } k \text {   records } l \text {   when } j \text {   is   correct }}{\text { number   of   patients   seen   by   observer } k \text {   where } j \text {   is   correct }}. \tag {1.1}
\]

However, in many cases the true response can never be ascertained.
Estimation of individual error-rates when the true response is not
available has been discussed by Dawid (1971) for a facet having just two
responses---yes or no, or 0 or 1. In the next section we extend the
analysis proposed by Dawid to a facet having several possible responses.
Section 3 mentions some of the computational problems and the practical
limitations of this method of analysis. Section 4 gives an example.

\section{2. MAXIMUM LIKELIHOOD
ESTIMATION}\label{maximum-likelihood-estimation}

Consider an experiment where K clinicians, indexed \(k = 1, \ldots, K\)
, ask a single question of I patients indexed \(i = 1, \ldots, I\) . Not
all clinicians need see every patient and a clinician may question the
same patient more than once. The answer received is one of J possible
replies. Let \(n_{il}^{(k)}\) be the number of times clinician k gets
responses l from patient i. The object of the experiment is to measure
the individual error-rates \(\pi_{jl}^{(k)}\) ( \(j = 1, \ldots, J\) ;
\(l = 1, \ldots, J\) ; \(k = 1, \ldots, K\) ). It is assumed that

\begin{enumerate}
\def\labelenumi{(\roman{enumi})}
\tightlist
\item
  The responses given by a single patient to successive clinicians or to
  repeated questioning by a single clinician are independent, given the
  true response. Further, all patients respond independently.\\
\item
  There is no patient-by-clinician interaction. For example, one
  clinician does not obtain a more helpful response from a patient than
  is given to other clinicians.
\end{enumerate}

The restrictive nature of these assumptions should be carefully noted in
any application of the following analysis. In particular, in many
studies there may be variables which intervene between the true response
and the elicited responses and which, for each patient, are common to
all clinicians. In the context of history-taking, such a variable might
be the patient's vague memory of the true response. The assumption (i)
above might more reasonably require independence of the responses given
the true response and the intervening variables. However, assumptions
(i) and (ii) are at present necessary if estimates of error rates are to
be obtained.

\section{Case 1. True responses are
available}\label{case-1.-true-responses-are-available}

Let \(\{T_{ij}: j=1,\ldots,J\}\) be a set of indicator variables for
patient i. If response q is the true response for this patient then
\(T_{iq}=1\) and \(T_{ij}=0\) ( \(j\neq q\) ). Further, the patients in
the experiment are considered to be a random sample from some
population, where the probability that a patient drawn at random has
true response j is \(p_{j}\) ( \(j=1,\ldots,J\) ). Typically these
probabilities are unknown.

Consider a single patient i and clinician k. Then, if q were the true
response, the numbers of responses of each type actually obtained would
be distributed according to a multinomial distribution and the
likelihood would be proportional to

\[
\prod_ {l = 1} ^ {J} \left(\pi_ {q l} ^ {(k)}\right) ^ {n _ {i l} ^ {(k)}}.
\]

As all clinicians elicit responses independently the likelihood for the
responses of patient i when \(T_{iq} = 1\) is

\[
\prod_ {k = 1} ^ {K} \prod_ {l = 1} ^ {J} \left(\pi_ {q l} ^ {(k)}\right) ^ {n _ {i l} ^ {(k)}}
\]

and unconditionally

\[
\prod_ {j = 1} ^ {J} \left\{p _ {j} \prod_ {k = 1} ^ {K} \prod_ {l = 1} ^ {J} \left(\pi_ {j l} ^ {(k)}\right) ^ {n _ {i l} ^ {(k)}} \right\} ^ {T _ {i j}}. \tag {2.1}
\]

That is, (2.1) is proportional to the probability of all the data (true
class and responses) on patient i. It consists of a product of J terms,
J-1 of which equal 1 (as \(T_{ij}=0, j\neq q\) ) and one term of the
form

\[
p (\text { responses   obtained } \mid T _ {i q} = 1) p (T _ {i q} = 1).
\]

As the data from all patients are assumed to be independent, the
likelihood for the full data is

\[
\prod_ {i = 1} ^ {I} \prod_ {j = 1} ^ {J} \left\{p _ {j} \prod_ {k = 1} ^ {K} \prod_ {l = 1} ^ {J} \left(\pi_ {j l} ^ {(k)}\right) ^ {n _ {i l} ^ {(k)}} \right\} ^ {T _ {i j}}. \tag {2.2}
\]

In (2.2) the quantities \(n_{ii}^{(k)}\) , \(T_{ij}\) and possibly
\(p_{j}\) are all known. The maximum likelihood estimates can be
calculated analytically, and we obtain estimators

\[
\hat {\pi} _ {j l} ^ {(k)} = \sum_ {i} T _ {i j} n _ {i l} ^ {(k)} / \sum_ {l} \sum_ {i} T _ {i j} n _ {i l} ^ {(k)}. \tag {2.3}
\]

The interpretation of (2.3) is simply equation (1.1). When the
probabilities \(p_j\) ( \(j = 1, \dots, J\) ) are unknown these can also
be estimated:

\[
\hat {p} _ {j} = \sum_ {i} T _ {i j} / I. \tag {2.4}
\]

We note, at this point, that when the individual error-rates
\(\{\pi_{jl}^{(kl)}\}\) and the marginal probabilities \(\{p_j\}\) are
known but the true class of a given patient, \(i\) , is unknown, Bayes'
Theorem may be used

to obtain estimates of the indicator variables \(T_{ij}\) (
\(j = 1, \dots, J\) ). A priori, \(p(T_{ij} = 1) = p_j\) . If data are
then collected from the patients and the counts \(n_{ii}^{(k)}\) are
obtained, by Bayes' Theorem

\[
p (T _ {i j} = 1 \mid \text { data }) \propto p (\text { data } \mid T _ {i j} = 1) p (T _ {i j} = 1)
\]

\[
\propto \prod_ {k = 1} ^ {K} \prod_ {l = 1} ^ {J} \left(\pi_ {j l} ^ {(k)}\right) ^ {n _ {i l} ^ {(k)}} p _ {j},
\]

where all terms not involving j are absorbed into the proportionality
sign. Thus

\[
p \left(T _ {i j} = 1 \mid \text { data }\right) = \prod_ {k = 1} ^ {K} \prod_ {l = 1} ^ {J} \left(\pi_ {j l} ^ {(k)}\right) ^ {n _ {u} ^ {(k)}} p _ {j} / \sum_ {q = 1} ^ {J} \prod_ {k = 1} ^ {K} \prod_ {l = 1} ^ {J} \left(\pi_ {q l} ^ {(k)}\right) ^ {n _ {u} ^ {(k)}} p _ {q}. \tag {2.5}
\]

\section{Case 2. True responses are not
available}\label{case-2.-true-responses-are-not-available}

In this case the probability for the data observed on a single patient
remains unchanged, given q is the true response. But as we do not know
which response is the true response, unconditionally

\[
p (\text { data   on   patient } i) \propto \sum_ {j = 1} ^ {J} p _ {j} \prod_ {k = 1} ^ {K} \prod_ {l = 1} ^ {J} \left(\pi_ {j l} ^ {(k)}\right) ^ {n _ {l l} ^ {(k)}}. \tag {2.6}
\]

Compare equations (2.6) and (2.1). In (2.1), as the \(T_{ij}\) (
\(j=1,\ldots,J\) ) were known, the distribution for the data on patient
i was essentially multinomial. Equation (2.6) is a mixture of such
multinomial distributions, the weights being the marginal probabilities
\(p_{j}\) ( \(j=1,\ldots,J\) ). The likelihood for full data is
therefore

\[
\prod_ {i = 1} ^ {I} \left(\sum_ {j = 1} ^ {J} p _ {j} \prod_ {k = 1} ^ {K} \prod_ {l = 1} ^ {J} \left(\pi_ {j l} ^ {(k)}\right) ^ {n _ {l l} ^ {(k)}}\right). \tag {2.7}
\]

Calculation of the maximum likelihood estimates for the \(p\) 's and
\(\pi\) 's in (2.7) is a much more difficult task. Dawid (1971)
discusses the problems in detail but explicit expressions comparable to
(2.3) and (2.4) cannot generally be obtained and numerical methods for
the solution of simultaneous equations or function-maximization have to
be employed.

In the past, the numerical methods available have not been satisfactory
because of the large number of unknown parameters involved. However,
Dempster et al.~(1977) describe a numerical method of maximum likelihood
estimation which is ideally suited to this particular problem. In
certain circumstances missing data preclude the straightforward maximum
likelihood estimation of the parameters of interest. However, if these
parameters are known the missing data can be estimated. An iterative
procedure is therefore proposed.

\begin{enumerate}
\def\labelenumi{(\roman{enumi})}
\tightlist
\item
  Obtain some initial estimates of the missing data.\\
\item
  Calculate the maximum likelihood estimates for the quantities of
  interest as if the missing data had been found.\\
\item
  Now calculate new estimates of the missing data.\\
\item
  Repeat steps (ii) and (iii) until both the maximum likelihood
  estimates and the missing data estimates converge.
\end{enumerate}

This procedure is known as the EM algorithm as each iteration consists
of an Expectation (of missing data) step and a Maximization (maximum
likelihood estimation) step. Dempster et al.~give conditions under which
the estimates obtained are those which should have been obtained if the
quantities of interest had been estimated directly by maximum likelihood
from the more complicated model which takes account of the missing data.

If, in the problem at hand, the indicator variables
\(\{T_{ij}; i=1,\ldots,I, j=1,\ldots,J\}\) are treated as missing data
then the conditions of the EM algorithm are satisfied. We therefore
proceed as follows:

\begin{enumerate}
\def\labelenumi{(\roman{enumi})}
\item
  Take initial estimates of the \(T\) 's.
\item
  Use equations (2.3) and (2.4) to obtain estimates of the \(p\) 's and
  \(\pi\) 's.\\
\item
  Use equation (2.5) and the estimates of the \(p\) 's and \(\pi\) 's to
  calculate new estimates of the \(T\) 's.\\
\item
  Repeat steps (ii) and (iii) until the results converge.
\end{enumerate}

In step (iii) the estimate of \(T_{ij}\) is
\(E(T_{ij} \mid \text{data}) = p(T_{ij} = 1 \mid \text{data})\) . That
is, the estimate is expressed as a probability that the true response
for patient \(i\) is \(j\) . Such probabilities may also be used as
initial input in step (i).

The final estimates of the \(p\) 's and \(\pi\) 's are those values
which maximize (2.7). The final estimates of the \(T\) 's are the
consensus probabilities for each patient from which the true response
for each patient can be assessed.

\section{3. DISCUSSION}\label{discussion}

A number of comments relating to the execution of the algorithm and the
interpretation of the results should be made at this point. Clearly, in
any application, the concept of a true response must be meaningful,
either as a hypothetically observable quantity or as the consensus
agreement of qualified medical opinion given complete information on the
patient. Even so, when the true responses are not known, there is no
unique way in which observer error-rates can be estimated as there is no
way in which the observations recorded by clinicians can be judged to be
sensible or even relevant to the true state of the patient. This is
reflected in the model, which has the structure of a latent class model
(Lazarsfeld and Henry, 1968), and thus suffers from the lack of
identifiability characteristic of factor-analytic models. Specifically,
the likelihood (2.7) remains unchanged for any relabelling of the j
index, and so it follows from the symmetry of the equation that there
are J! sets of estimates which correspond to the global maximum. Some
will be more sensible than others, and it should be sufficient to assume
that the correct estimates are those where
\(\pi_{ij}^{(k)} > \pi_{jl}^{(k)} (j \neq l)\) for most k and j.

Closely allied to the issue of identifiability is the question of
initial estimates. One possibility is to assign
\(T_{ij} = 1 / J (i = 1, \dots, I, j = 1, \dots, J)\) . However, this
action should be avoided as the initial estimates of the \(p\) 's and
\(\pi\) 's correspond exactly to the centre of symmetry in (2.7)---a
saddle point---and, as noted by Dempster et al., the EM algorithm cannot
converge from such a point. A second possibility is to use the data to
calculate initial estimates. For example, one could use

\[
\hat {T} _ {i j} = \sum_ {k} n _ {i l} ^ {(k)} / \sum_ {k} \sum_ {l} n _ {i l} ^ {(k)} \tag {3.1}
\]

as starting values.

In practice, it is advisable to repeat the algorithm for several
different sets of starting values. The EM algorithm is only guaranteed
to converge to a local maximum and in situations where relatively few
data are used to estimate a large number of parameters one must usually
be content with choosing from a set of estimates corresponding to
different local maxima on the basis of the magnitude of the likelihood.
In our experience the starting values defined by (3.1) have been
particularly useful in locating the best local maximum, if not the
global maximum of interest. A maximum likelihood algorithm closely
related to that above is given by Goodman (1974) who similarly advises
on the use of several different starting values.

Little is known at present about the accuracy of the estimates. It is
suspected that the standard errors of each of the \(\pi\) 's and the p's
are large and that these estimates are highly correlated. However, the
large number of parameters makes the dispersion matrix difficult to
obtain and, given that the maximum likelihood estimates will almost
invariably lie on the boundary of the parameter-space, of limited
usefulness. One approach has been to observe the stability of the
estimates to the removal of individual patients or individual observers
from the data. Our conclusion is that the Example below probably
represents the smallest

experiment which will yield point estimates of any value unless the raw
data indicate a very high level of agreement.

One consequence of probability estimates which are either zero or one is
that unrealistically accurate estimates of the T's are obtained.
However, inspection of the final estimates of the T's during repeated
computation of the error-rates can lead to a clear appreciation of the
``true'' class of each patient in the experiment. While the vast
majority of the final estimates of the T's remain virtually unaltered as
different local maxima are discovered, some T estimates change markedly
indicating the ``reclassification'' of one or two patients. A similar
feature is noticed when the algorithm is repeated following removal of
cases. Thus it is very easy to establish which patients are well
classified and for which patients the consensus is more uncertain.

Dempster et al.~prove that the algorithm exhibits first-order
convergence. The significance of this result in our context depends
largely on the proportion of cases in the data set where a consensus is
not immediately obvious. If this proportion is substantial---say 30 per
cent or more---a large number of iterations is usually required and
repeated use of the algorithm becomes expensive. Nelder (1977) suggests
that a substantial improvement is possible by exploring ahead along the
apparent direction of convergence. This procedure has been used with
good effect.

Finally, it must be emphasized that the method of maximum likelihood
estimation has few advantages other than tractability. For a model as
highly parameterized as this, having \((J-1)(JK+1)\) parameters, it
would be naive to expect any of the theoretical large sample optimality
properties to hold. It would be desirable to develop other estimation
methods free of the general criticisms and caveats given above and to
allow for prior opinions about the relative accuracies of the various
questions. One possibility is to express the \(\pi\) 's as functions of
a much smaller number of parameters and estimate these by maximum
likelihood. We hope to present this method in a future paper.

\section{4. AN EXAMPLE}\label{an-example}

In the pre-operative assessment of a patient an anaesthetist must decide
whether a patient is fit enough to undergo a general anaesthetic. In the
trial reported here a standard form was completed on 45 patients by an
independent party and contained information reflecting the

TABLE 1 Assessments of fitness for anaesthesia

Patient

Observer

Patient

Observer

Patient

Observer

1

2

3

4

5

1

2

3

4

5

1

2

3

4

5

1

111

1

1

1

1

16

111

2

1

1

1

31

111

1

1

1

1

2

333

4

3

3

4

17

111

1

1

1

1

32

333

3

2

3

3

3

112

2

1

2

2

18

111

1

1

1

1

33

111

1

1

1

1

4

222

3

1

2

1

19

222

2

2

2

1

34

222

2

2

2

2

5

222

3

2

2

2

20

222

1

3

2

2

35

222

3

2

3

2

6

222

3

3

2

2

21

222

2

2

2

2

36

433

4

3

4

3

7

122

2

1

1

1

22

222

2

2

2

1

37

221

2

2

3

2

8

333

3

4

3

3

23

222

3

2

2

2

38

232

3

2

3

3

9

222

2

2

2

3

24

221

2

2

2

2

39

333

3

4

3

2

10

232

2

2

2

3

25

111

1

1

1

1

40

111

1

1

1

1

11

444

4

4

4

4

26

111

1

1

1

1

41

111

1

1

1

1

12

222

3

3

4

3

27

232

2

2

2

2

42

121

2

1

1

1

13

111

1

1

1

1

28

111

1

1

1

1

43

232

2

2

2

2

14

222

3

2

1

2

29

111

1

1

1

1

44

121

1

1

1

1

15

121

1

1

1

1

30

112

1

1

2

1

45

222

2

2

2

2

patient's state of health. These forms were then independently assessed
by five anaesthetists who classified each patient on a 1 to 4 scale,
these categories having previously been defined. Anaesthetist 1 assessed
the forms a total of three times, each time separated by some weeks. The
data are shown in Table 1. Table 2 gives the estimates of the marginal
probabilities and the individual error-rates. A range of starting values
were explored although those described by (3.1) yielded the best maximum
in this instance. Table 3 gives the matrices of probabilities
\((p_{j}\pi_{jl}^{(k)})\) for each observer. In many circumstances these
incidence-of-error matrices are of more immediate interest than the
error-rates. The sum of the diagonal elements of such a matrix yields an
estimate of the probability of a correct allocation by an observer and,
when the feature under consideration is measured on an ordinal scale as
in this example, the sum

TABLE 2 Maximum likelihood estimates

Marginal probabilities

Category:

1

2

3

4

·40

·42

·11

·07

Error-rates

OBSERVER 1

Observed response:

1

2

3

4

True response

1

·89

·11

·0

·0

2

·07

·88

·05

·0

3

·0

·34

·66

·0

4

·0

·0

·56

·44

OBSERVER 2

Observed response:

1

2

3

4

True response

1

·78

·22

·0

·0

2

·06

·84

·10

·0

3

·0

·0

1·0

·0

4

·0

·0

·0

1·0

OBSERVER 3

Observed response:

1

2

3

4

True response

1

1·0

·0

·0

·0

2

·12

·79

·09

·0

3

·0

·40

·20

·4

4

·0

·0

·67

·33

OBSERVER 4

Observed response:

1

2

3

4

True response

1

·94

·06

·0

·0

2

·05

·84

·11

·0

3

·0

·0

·80

·20

4

·0

·0

·33

·67

OBSERVER 5

Observed response:

1

2

3

4

True response

1

1·0

·0

·0

·0

2

·16

·74

·10

·0

3

·0

·21

·79

·0

4

·0

·0

·33

·67

TABLE 3 Incidence-of-error probabilities

Observed response:

OBSERVER 1

1

2

3

4

True response 1

·36

·04

·0

·0

2

·03

·37

·02

·0

3

·0

·04

·07

·0

4

·0

·0

·04

·03

Observed response:

OBSERVER 2

1

2

3

4

True response 1

·31

·09

·0

·0

2

·02

·27

·13

·0

3

·0

·0

·11

·0

4

·0

·0

·0

·07

Observed response:

OBSERVER 3

1

2

3

4

True response 1

·40

·0

·0

·0

2

·05

·33

·04

·0

3

·0

·045

·02

·045

4

·0

·0

·05

·02

Observed response:

OBSERVER 4

1

2

3

4

True response 1

·38

·02

·0

·0

2

·02

·36

·04

·0

3

·0

·0

·09

·02

4

·0

·0

·02

·05

Observed response:

OBSERVER 5

1

2

3

4

True response 1

·40

·0

·0

·0

2

·07

·31

·04

·0

3

·0

·02

·09

·0

4

·0

·0

·02

·05

of the elements below or above the diagonal reflect the observer's
optimism or pessimism relative to the other observers in the trial.

Table 4 gives the estimated probabilities for the T's for each patient.
For most patients the posterior probability is 1·0 for a single
response. In these cases the consensus appears obvious, but as mentioned
above these results could be viewed with some suspicion. In fact,

TABLE 4\\
Final estimates of indicator variables for each patient†

Patient

Category

Patient

Category

1

2

3

4

1

2

3

4

1

1·0

0·0

0·0

0·0

24

0·0

1·0

0·0

0·0

2

0·0

0·0

0·0

1·0

25

1·0

0·0

0·0

0·0

3

0·0

1·0

0·0

0·0

26

1·0

0·0

0·0

0·0

4

0·0

1·0

0·0

0·0

27

0·0

1·0

0·0

0·0

5

0·0

1·0

0·0

0·0

28

1·0

0·0

0·0

0·0

6

0·0

1·0

0·0

0·0

29

1·0

0·0

0·0

0·0

7

0·986

0·014

0·0

0·0

30

0·999

0·001

0·0

0·0

8

0·0

0·0

1·0

0·0

31

1·0

0·0

0·0

0·0

9

0·0

1·0

0·0

0·0

32

0·0

0·0

1·0

0·0

10

0·0

1·0

0·0

0·0

33

1·0

0·0

0·0

0·0

11

0·0

0·0

0·0

1·0

34

0·0

1·0

0·0

0·0

12

0·0

0·0

1·0

0·0

35

0·0

0·948

0·052

0·0

13

1·0

0·0

0·0

0·0

36

0·0

0·0

0·0

1·0

14

0·0

1·0

0·0

0·0

37

0·0

1·0

0·0

0·0

15

1·0

0·0

0·0

0·0

38

0·0

0·021

0·979

0·0

16

1·0

0·0

0·0

0·0

39

0·0

0·0

1·0

0·0

17

1·0

0·0

0·0

0·0

40

1·0

0·0

0·0

0·0

18

1·0

0·0

0·0

0·0

41

1·0

0·0

0·0

0·0

19

0·0

1·0

0·0

0·0

42

1·0

0·0

0·0

0·0

20

0·0

1·0

0·0

0·0

43

0·0

1·0

0·0

0·0

21

0·0

1·0

0·0

0·0

44

1·0

0·0

0·0

0·0

22

0·0

1·0

0·0

0·0

45

0·0

1·0

0·0

0·0

23

0·0

1·0

0·0

0·0

† Probabilities of 1·0 and 0·0 are correct to three decimal places.

repeated trials suggest that patients numbered 2, 12, 14 and 38 are the
only cases which are sensitive to small changes in error-rate estimates,
and thus a definite statement of true class is possible in many cases
where consensus is not apparent in the raw data.

\section{ACKNOWLEDGEMENT}\label{acknowledgement}

We are grateful to Dr M. E. Wilson of the Bristol Royal Infirmary for
permission to use the data in Section 4.

\section{REFERENCES}\label{references}

DAVID, A. P. (1971). Estimation of error-rates in history taking.
Preliminary analysis. Paper presented to the Royal College of Physicians
Computer Workshop, November 1971.\\
DEMPSTER, A. P., LAIRD, N. M. and RUBIN, D. B. (1977). Maximum
likelihood from incomplete data via the EM algorithm (with Discussion).
J. R. Statist. Soc. B, 39, 1--38.\\
GOOD, I. J. and CARD, W. I. (1971). The diagnostic process with special
reference to errors. Meth. Inform. Med., 10, 176--188.\\
GOODMAN, L. A. (1974). Exploratory latent structure analysis using both
identifiable and unidentifiable models. Biometrika, 61, 215--231.

LANDIS, J. R. and KOCH, G. G. (1975a). A review of statistical methods
in the analysis of data arising from observer reliability studies (Part
I). Statist. Neerland., 29, 101--123.\\
(1975b). A review of statistical methods in the analysis of data arising
from observer reliability studies (Part II). Statist. Neerland., 29,
151--161.\\
--- (1977). The measurement of observer agreement for categorical data.
Biometrics, 33, 159--174.\\
LAZARSFELD, P. F., and HENRY, N. W. (1968). Latent Structure Analysis.
Boston: Houghton Mifflin. NELDER, J. A. (1977). In discussion of
Dempster et al.~(1977).

\end{document}
